\documentclass[11pt,a4paper]{article}
\input{../../sablony/pripravy/priprava.sty}
\usepackage{tikz,pgfplots}
\pgfplotsset{compat=1.18}
\nastavPripravu{1. ročník}{Kvadratické rovnice a nerovnice}{109}{Kvadratické rovnice: volba postupu}
\begin{document}
\blok{Kvadratické rovnice: volba postupu}
Souhrnné procvičení neúplných a úplných rovnic, nulového součinu a Viètových vztahů.
\blok{Metoda a návaznost na graf}
\begin{itemize}
\item Žák před výpočtem řekne, proč zvolil odmocnění, vytknutí, rozklad či diskriminant.
\item Při řešení kontroluj počet kořenů a souvislost s grafem.
\end{itemize}
\blok{Společně řešené příklady}
\begin{prikladReseniBox}[title={Příklad 1}]
$3x^{2} - 48=0$
\tcblower
\[
\begin{aligned}
x^2&=16\\|x|&=4\\[2pt]
K&=\{-4;4\}
\end{aligned}
\]
\end{prikladReseniBox}
\begin{prikladReseniBox}[title={Příklad 2}]
$2x^{2} - 10x=0$
\tcblower
\[
\begin{aligned}
2x^{2} - 10x&=x(2x - 10)=0\\[2pt]
K&=\{0;5\}
\end{aligned}
\]
\end{prikladReseniBox}
\begin{prikladReseniBox}[title={Příklad 3}]
$x^{2} - x - 6=0$
\tcblower
\[
\begin{aligned}
&\left(x - 3\right) \cdot \left(x + 2\right)=0\\[2pt]
x&=3\ \vee\ x=-2\\[2pt]
K&=\{-2;3\}
\end{aligned}
\]
\end{prikladReseniBox}
\blok{Samostatné procvičování}
\begin{zadaniReseniBox}[title={Úloha 1}]
$2x^{2} - 7x + 3=0$
\tcblower
\[
\begin{aligned}
D&=49-4\cdot(2)\cdot(3)=25\\[2pt]
x_1&=\frac{1}{2}\\[2pt]
x_2&=3\\[2pt]
K&=\{\frac{1}{2};3\}
\end{aligned}
\]
\end{zadaniReseniBox}
\begin{zadaniReseniBox}[title={Úloha 2}]
$x^{2} + 6x + 9=0$
\tcblower
\[
\begin{aligned}
D&=36-4\cdot(1)\cdot(9)=0\\[2pt]
x&=-3\\[2pt]
K&=\{-3\}
\end{aligned}
\]
\end{zadaniReseniBox}
\blok{Volitelná rezerva}
\poznamka{Studenti si vymyslí obdobný příklad se změnou jednoho koeficientu a vysvětlí, jak se změní počet kořenů či řešení.}
\blok{Rozšiřující příklady ze sbírky}
\poznamka{Pravidlo nulového součinu funguje i pro součin více polynomů. Druhá úloha má jeden kvadratický činitel bez reálných nul.}

\begin{zadaniReseniBox}[title={Rozšiřující příklad}]
Řešte v $\mathbb R$: $- \left(4x^{2} - 9\right) \cdot \left(x^{2} - 6x + 9\right)=0$.
\tcblower
\[
\begin{aligned}- \left(x - 3\right)^{2} \cdot \left(2x - 3\right) \cdot \left(2x + 3\right)&=0\\K&=\{- \frac{3}{2};\frac{3}{2};3\}\end{aligned}
\]
\end{zadaniReseniBox}

\begin{zadaniReseniBox}[title={Rozšiřující příklad}]
Řešte v $\mathbb R$: $- 4x^{4} + 175x^{2} + 1250=0$.
\tcblower
\[
\begin{aligned}- 4x^{4} + 175x^{2} + 1250&=0\\K&=\{- 5 \cdot \sqrt{2};5 \cdot \sqrt{2}\}\end{aligned}
\]
\end{zadaniReseniBox}

\blok{Volitelná výzva -- různé metody}
\begin{zadaniReseniBox}[title={Rozšiřující příklad}]
$x^2=7x-10$.
\tcblower
\[\begin{aligned}
x^2-7x+10&=0\\(x-2)(x-5)&=0\\K&=\{2;5\}
\end{aligned}\]
\end{zadaniReseniBox}
\blok{Domácí úkol}
\begin{zadaniReseniBox}[title={Domácí příprava}]
\begin{enumerate}[label=\alph*),itemsep=2pt]
\item $4x^{2} - 25=0$
\item $x^{2} + 16=0$
\item $x^{2} - 9x + 14=0$
\item $2x^{2} + x - 6=0$
\item $4x^{2} + 12x + 9=0$
\item $3x^{2} - 5x + 1=0$
\end{enumerate}
\tcblower
\begin{enumerate}[label=\alph*),itemsep=2pt]
\item
\[\begin{aligned}
x^2&=\frac{25}{4}\\|x|&=\frac52\\[2pt]
K&=\{- \frac{5}{2};\frac{5}{2}\}
\end{aligned}\]
\item
\[\begin{aligned}
x^2&=-16\\[2pt]
K&=\varnothing
\end{aligned}\]
\item
\[\begin{aligned}
&\left(x - 7\right) \cdot \left(x - 2\right)=0\\[2pt]
x&=7\ \vee\ x=2\\[2pt]
K&=\{2;7\}
\end{aligned}\]
\item
\[\begin{aligned}
D&=1-4\cdot(2)\cdot(-6)=49\\[2pt]
x_1&=-2\\[2pt]
x_2&=\frac{3}{2}\\[2pt]
K&=\{-2;\frac{3}{2}\}
\end{aligned}\]
\item
\[\begin{aligned}
D&=144-4\cdot(4)\cdot(9)=0\\[2pt]
x&=- \frac{3}{2}\\[2pt]
K&=\{- \frac{3}{2}\}
\end{aligned}\]
\item
\[\begin{aligned}
D&=25-4\cdot(3)\cdot(1)=13\\[2pt]
x_1&=\frac{5}{6} - \frac{\sqrt{13}}{6}\\[2pt]
x_2&=\frac{\sqrt{13}}{6} + \frac{5}{6}\\[2pt]
K&=\{\frac{5}{6} - \frac{\sqrt{13}}{6};\frac{\sqrt{13}}{6} + \frac{5}{6}\}
\end{aligned}\]
\end{enumerate}
\end{zadaniReseniBox}
\poznamka{Orientační čas: úvod 5 min, společné příklady 18 min, samostatně 15 min, kontrola a shrnutí 7 min. Rezerva je dobrovolná. Číselné parametry modelových aplikací nejsou skutečně naměřená data.}
\end{document}
